\documentclass[11pt]{article}
\usepackage[a4paper,margin=1in]{geometry}
\usepackage{fontspec}
\usepackage[boldfont=false]{xeCJK}
\setCJKmainfont{FandolSong-Regular.otf}
\usepackage{amsmath}
\usepackage{amssymb}
\usepackage{array}
\usepackage{booktabs}
\usepackage{graphicx}
\usepackage{adjustbox}
\usepackage{enumitem}
\setlist[itemize]{topsep=2pt,partopsep=0pt,parsep=0pt,itemsep=2pt,leftmargin=1.4em}
\setlist[enumerate]{topsep=2pt,partopsep=0pt,parsep=0pt,itemsep=2pt,leftmargin=1.6em}
\usepackage{parskip}
\usepackage{fvextra}
\fvset{breaklines=true,breakanywhere=true,fontsize=\small}
\setlength{\parindent}{0pt}

\begin{document}

\begin{center}
  {\LARGE\bfseries Chemical bonding}\\[0.35em]
  {\large A Level Chemistry}
\end{center}
\vspace{0.6em}

\section*{Electronegativity}

\textbf{Electronegativity} 电负性 is the power of an atom to attract the \textbf{electrons} 电子 in a bond towards itself.

Three factors decide how electronegative an atom is:

\begin{itemize}
\item \textbf{nuclear charge} 核电荷: more protons pull the bonding electrons more strongly.

\item \textbf{atomic radius} 原子半径: the closer the bond is to the nucleus, the stronger the pull.

\item \textbf{shielding} 屏蔽 by inner shells and sub-shells: more inner electrons weaken the pull on the bonding electrons.

\end{itemize}

So electronegativity \textbf{rises} across a period (more nuclear charge, smaller radius) and \textbf{falls} down a group (larger radius, more shielding). Fluorine is the most electronegative element.

You can use the difference in \textbf{Pauling electronegativity} 鲍林电负性 values to predict the bond type. A large difference gives an \textbf{ionic} bond; a small difference gives a \textbf{covalent} bond.

\section*{Ionic bonding}

\textbf{Ionic bonding} 离子键 is the \textbf{electrostatic attraction} 静电引力 between oppositely charged \textbf{ions} 离子 — positive \textbf{cations} 阳离子 and negative \textbf{anions} 阴离子.

It forms when a metal gives electrons to a non-metal. Good examples are sodium chloride ($\text{NaCl}$), magnesium oxide ($\text{MgO}$) and calcium fluoride ($\text{CaF}_2$). The ions pack into a regular giant \textbf{lattice} 晶格, held together by the attraction in every direction.

\section*{Metallic bonding}

\textbf{Metallic bonding} 金属键 is the electrostatic attraction between positive metal ions and a "sea" of \textbf{delocalised electrons} 离域电子.

The outer electrons are free to move through the whole metal. This explains why metals conduct electricity and are strong.

\section*{Covalent and coordinate bonding}

\textbf{Covalent bonding} 共价键 is the electrostatic attraction between the nuclei of two atoms and a shared pair of electrons.

Simple molecules with covalent bonds include $\text{H}_2$, $\text{O}_2$, $\text{N}_2$, $\text{Cl}_2$, $\text{HCl}$, $\text{CO}_2$, $\text{NH}_3$, $\text{CH}_4$, $\text{C}_2\text{H}_6$ and $\text{C}_2\text{H}_4$. A double bond shares two pairs; a triple bond (as in $\text{N}_2$) shares three pairs.

Atoms in Period 3 and below can \textbf{expand the octet} 扩展八隅体 — hold more than eight electrons in their outer shell. Examples are $\text{SO}_2$, $\text{PCl}_5$ and $\text{SF}_6$.

A \textbf{coordinate bond} 配位键 (also called a dative covalent bond) is a covalent bond where both shared electrons come from the same atom. For example, when ammonia and hydrogen chloride gases meet, the lone pair on the nitrogen forms a coordinate bond to $\text{H}^+$, making the ammonium ion $\text{NH}_4^+$. Coordinate bonds also join the two halves of the $\text{Al}_2\text{Cl}_6$ molecule.

\subsection*{Sigma and pi bonds}

Covalent bonds form when \textbf{orbitals} 轨道 overlap:

\begin{itemize}
\item a \textbf{sigma bond} σ键 forms by the direct, head-on \textbf{overlap} 重叠 of orbitals between the two atoms.

\item a \textbf{pi bond} π键 forms by the sideways overlap of two p orbitals, above and below the sigma bond.

\end{itemize}

A single bond is one sigma bond. A double bond (as in $\text{C}_2\text{H}_4$) is one sigma plus one pi bond. A triple bond (as in $\text{N}_2$ and $\text{HCN}$) is one sigma plus two pi bonds.

\subsection*{Hybridisation}

\textbf{Hybridisation} 杂化 mixes orbitals in the same shell to make new, equal orbitals for bonding:

\begin{itemize}
\item $\text{sp}$: two equal orbitals, used in a linear molecule.

\item $\text{sp}^2$: three equal orbitals, used in a flat molecule like $\text{C}_2\text{H}_4$.

\item $\text{sp}^3$: four equal orbitals, used in $\text{CH}_4$.

\end{itemize}

\subsection*{Bond energy and bond length}

\begin{itemize}
\item \textbf{bond energy} 键能 is the energy needed to break one mole of a particular covalent bond in the gas state.

\item \textbf{bond length} 键长 is the distance between the centres of the two bonded atoms.

\end{itemize}

A shorter bond is usually stronger (higher bond energy). Triple bonds are shorter and stronger than double bonds, which are shorter and stronger than single bonds. Stronger bonds make a molecule harder to react.

\section*{Shapes of molecules}

To work out a shape, use \textbf{VSEPR theory} 价层电子对互斥理论: the pairs of electrons around the central atom push apart as far as possible, because like charges repel.

A \textbf{lone pair} 孤对电子 (not in a bond) pushes more strongly than a \textbf{bonding pair} 成键电子对. Each lone pair squeezes the \textbf{bond angle} 键角 by about $2.5°$.

\begin{center}
\adjustbox{max width=\linewidth}{%
\begin{tabular}{lll}
\toprule
\textbf{Molecule} & \textbf{Shape} & \textbf{Bond angle} \\
\midrule
$\text{CO}_2$ & \textbf{linear} 直线形 & $180°$ \\
$\text{BF}_3$ & \textbf{trigonal planar} 平面三角形 & $120°$ \\
$\text{CH}_4$ & \textbf{tetrahedral} 四面体形 & $109.5°$ \\
$\text{NH}_3$ & \textbf{pyramidal} 三角锥形 & $107°$ \\
$\text{H}_2\text{O}$ & \textbf{bent} 角形 & $104.5°$ \\
$\text{PF}_5$ & \textbf{trigonal bipyramidal} 三角双锥形 & $120°$ and $90°$ \\
$\text{SF}_6$ & \textbf{octahedral} 八面体形 & $90°$ \\
\bottomrule
\end{tabular}}
\end{center}

$\text{NH}_3$ has one lone pair and $\text{H}_2\text{O}$ has two, which is why their angles drop below the $109.5°$ of $\text{CH}_4$. You can predict the shapes of similar molecules and ions in the same way.

\section*{Intermolecular forces}

\textbf{Intermolecular forces} 分子间作用力 are the forces \emph{between} \textbf{molecules} 分子. They are much weaker than the ionic, covalent and metallic bonding \emph{inside} substances.

\subsection*{Bond polarity and dipoles}

When two atoms with different electronegativity share a bond, the electrons sit closer to the more electronegative atom. The bond then has a \textbf{polarity} 极性: one end is slightly negative ($\delta-$) and the other slightly positive ($\delta+$). This separation of charge is a \textbf{dipole} 偶极.

If the dipoles in a molecule do not cancel, the whole molecule has a \textbf{dipole moment} 偶极矩 and is polar. If they cancel by symmetry (as in $\text{CO}_2$), the molecule is non-polar.

\subsection*{Van der Waals' forces}

\textbf{Van der Waals' forces} 范德华力 is the general name for all intermolecular forces. There are two main types.

The first type is the instantaneous dipole–induced dipole force, also called the \textbf{London dispersion force} 伦敦色散力. Moving electrons make a brief \textbf{instantaneous dipole} 瞬时偶极, which then creates a matching \textbf{induced dipole} 诱导偶极 in a nearby molecule. These forces act between all molecules and get stronger when there are more electrons.

The second type is the permanent dipole–permanent dipole force. It acts between molecules that are always polar, because each one has a \textbf{permanent dipole} 永久偶极.

\subsection*{Hydrogen bonding}

\textbf{Hydrogen bonding} 氢键 is a strong, special case of permanent dipole forces. It forms when hydrogen is bonded to a very electronegative atom — nitrogen, oxygen or fluorine — and is attracted to a lone pair on an N, O or F atom in a neighbour. Look for N–H and O–H groups, as in ammonia and water.

Hydrogen bonding explains the strange behaviour of water:

\begin{itemize}
\item its high melting and \textbf{boiling point} 沸点, because many hydrogen bonds must be broken.

\item its high \textbf{surface tension} 表面张力.

\item ice is \textbf{less dense} than liquid water, because hydrogen bonds hold the molecules in an open, spread-out structure, so ice floats.

\end{itemize}

\section*{Dot-and-cross diagrams}

A \textbf{dot-and-cross diagram} 点叉图 shows the outer electrons of each atom, using dots for one atom and crosses for the other. This makes it clear where each bonding electron came from. You can draw them for ionic, covalent and coordinate bonding, including molecules with an expanded octet or an odd number of electrons.


\end{document}
